\documentclass[11pt,a4paper]{article}

\usepackage[margin=1.75cm]{geometry}
\usepackage{fontspec}
\usepackage{enumitem}
\usepackage{xurl}
\usepackage[hidelinks]{hyperref}
\usepackage{xcolor}
\usepackage{titlesec}
\usepackage{tabularx}

\setmainfont{Times New Roman}
\definecolor{heading}{HTML}{003366}
\definecolor{subtle}{HTML}{555555}

\pagestyle{empty}
\setlength{\parindent}{0pt}
\setlength{\parskip}{4pt}
\sloppy
\setlist[itemize]{leftmargin=1.25em, topsep=2pt, itemsep=1pt}

\titleformat{\section}{\large\bfseries\color{heading}}{}{0pt}{}[\titlerule]
\titleformat{\subsection}{\normalsize\bfseries\color{heading}}{}{0pt}{}

\newcommand{\doi}[1]{\href{https://doi.org/#1}{\nolinkurl{#1}}}
\newcommand{\pub}[6]{\item #1 (#2). #3. \textit{#4}. DOI: \doi{#5}. #6}

\begin{document}

{\LARGE \textbf{Wen Quan}}\\[-1pt]
\textcolor{subtle}{Academic CV for PhD Application in Computer Science and Technology}\\[3pt]
\href{mailto:wenquan@student.usm.my}{wenquan@student.usm.my} \,|\, 
\href{mailto:wenquan@hngm.edu.cn}{wenquan@hngm.edu.cn} \,|\, 
+86 156-4900-6328 \,|\, +60 12-938-6817\\
Academic homepage: \href{https://wenquan0816.github.io/schoolar-site/}{\nolinkurl{wenquan0816.github.io/schoolar-site}} \,|\, 
Portfolio: \href{https://wenquan0816.github.io/age-friendly-community-map/}{\nolinkurl{age-friendly-community-map}}\\
ORCID: \href{https://orcid.org/0009-0005-3975-8104}{0009-0005-3975-8104}

\section*{Application Objective}
Application for the \textbf{Doctor of Philosophy in Computer Science and Technology} at Xiamen University Malaysia, with a proposed research focus on \textbf{knowledge-guided vision-language models, lightweight spatial intelligence, and AI-assisted assessment of age-friendly indoor environments}.

\section*{Profile}
AI-assisted built environment researcher with a completed PhD in Interior Design from Universiti Sains Malaysia. My research connects age-friendly built environments, evidence-based design, public health resilience, spatial assessment, and applied AI methods. My published work has used LSTM-based model predictive control, natural language processing, SciBERT-assisted indicator extraction, contrastive learning, hyperbolic embedding, GANs, random forest regression, and multi-criteria assessment.

For the proposed PhD in Computer Science and Technology, my goal is to move from domain-specific AI applications toward computational methods for \textbf{knowledge-guided vision-language models and explainable spatial assessment systems} for age-friendly indoor environments.

\section*{Proposed PhD Research}
\textbf{Proposed title:} A Knowledge-Guided Vision-Language Framework for Age-Friendly Indoor Space Assessment and Design Recommendation.

\textbf{Research fit:} Artificial Intelligence, Human-Computer Interaction, Software Technology, multimodal AI, knowledge-based systems, and applied computer vision for built environment assessment.

\section*{Education}
\begin{itemize}
  \item \textbf{PhD in Interior Design}, School of Housing, Building and Planning, Universiti Sains Malaysia, Malaysia, 2021--2026. Status: completed / viva passed. Research area: age-friendly built environments, AI-assisted assessment, interior and community environment evaluation.
  \item \textbf{Master of Engineering in Architecture}, Henan University of Technology, China, 2014--2017. Specialization: Architectural Design and Theory.
  \item \textbf{Bachelor of Engineering in Architecture}, Anyang Normal University, China, 2009--2014. Dual Degree Program: Bachelor of Education.
\end{itemize}

\section*{Research Interests}
\begin{itemize}
  \item Vision-language models for spatial assessment and design recommendation.
  \item AI-assisted age-friendly indoor and community environment evaluation.
  \item Knowledge-guided risk assessment, explainable AI, and decision support.
  \item Human-computer interaction for design and built environment applications.
  \item Natural language processing for design guidelines, policies, and expert knowledge extraction.
\end{itemize}

\section*{Computational and AI Experience}
\begin{itemize}
  \item LSTM-based model predictive control for building energy flexibility.
  \item NLP and SciBERT-assisted framework development for community resilience and age-friendly conversion studies.
  \item Contrastive learning and hyperbolic embedding for public health emergency resilience assessment.
  \item GANs and random forest regression for assessment and prediction modelling.
  \item Meta-learning, attentive neural processes, 3D point-cloud features, ensemble learning, and SHAP-based explainable AI in manuscripts currently under review.
  \item Web-based interactive assessment demo using HTML, CSS, JavaScript, Leaflet, and ECharts.
  \item Multi-criteria assessment and expert validation: AHP, indicator weighting, framework development.
\end{itemize}

\section*{Selected Portfolio}
\textbf{China Age-Friendly Community Assessment Map}\\
Live demo: \href{https://wenquan0816.github.io/age-friendly-community-map/}{\nolinkurl{wenquan0816.github.io/age-friendly-community-map}}\\
Repository: \href{https://github.com/WENQUAN0816/age-friendly-community-map}{\nolinkurl{github.com/WENQUAN0816/age-friendly-community-map}}

This web-based demo visualizes multidimensional age-friendly community assessment results through an interactive map and dashboard. It demonstrates my ability to translate built environment research into computational scoring, spatial visualization, and user-facing prototype development.

\section*{Selected AI/CS Manuscripts Under Review}
\begin{enumerate}
  \item Quan Wen; Yanting Wu; Mazran Ismail; Muhammad Hafeez Abdul Nasir. \textbf{Meta-Learning Framework for Healthy Aging Assessment: Attentive Neural Processes with Cross-Region Generalization within CHARLS}. \textit{Neural Networks}. Status: R1 revised submission under processing, NEUNET-D-26-01296R1; first author and corresponding author.
  \item Quan Wen; Mazran Ismail; Muhammad Hafeez Abdul Nasir; Nooriati Binti Taib; Jestin Bin Nordin. \textbf{Spatial fall risk prediction in elderly residential environments using 3D point cloud-derived features and ensemble learning with SHAP interpretability}. \textit{Egyptian Informatics Journal}. Status: R1 revised submission under processing, EGIJ-D-26-00407R1; first author and corresponding author.
\end{enumerate}

\section*{Publications}
\subsection*{WoS SCIE/SSCI Publications}
\begin{enumerate}
  \pub{Quan Wen; Mazran Ismail; Muhammad Hafeez Abdul Nasir}{2026}{A LSTM-based model predictive control method for unlocking the potential of building energy flexibility in Malaysian commercial buildings}{Scientific Reports}{10.1038/s41598-026-53501-8}{WoS SCIE Q1; first author.}
  \pub{Quan Wen; Mazran Ismail; Muhammad Hafeez Abdul Nasir; Yanting Wu; Chaojiang Hu}{2026}{AFLE-HIV: developing an age-friendly living environment assessment framework for elderly HIV-infected individuals in China}{BMC Public Health}{10.1186/s12889-026-27748-9}{WoS SCIE Q1; first author.}
  \pub{Quan Wen; Mazran Ismail; Muhammad Hafeez Abdul Nasir}{2026}{Transforming vacant rural school buildings into age-friendly facilities under China's aging and low-fertility context: A natural language processing-based evaluation framework}{BMC Geriatrics}{10.1186/s12877-025-06772-1}{WoS SCIE Q1; first author.}
  \pub{Quan Wen; Mazran Ismail; Muhammad Hafeez Abdul Nasir}{2025}{A community public health emergency resilience assessment framework based on contrastive learning and hyperbolic embedding}{Frontiers in Public Health}{10.3389/fpubh.2025.1651331}{WoS SSCI Q1; first author.}
  \pub{Yanli Han; Quan Wen}{2026}{Emotion Recognition in Literary Texts via an Attention-Guided Dual-Stream Gated Fusion Network}{IEEE Access}{10.1109/ACCESS.2026.3686065}{WoS SCIE; second author; corresponding author.}
  \pub{Wenjing Xu; Mazran Ismail; Syafizal Shahruddin; Wen Quan; Yingrui Li}{2025}{Exploring tourists' intentions to adopt augmented reality in cultural heritage museums: Insights from a modified technology acceptance model}{SAGE Open}{10.1177/21582440251339936}{WoS SSCI Q1; co-author.}
\end{enumerate}

\subsection*{WoS ESCI and Scopus-Indexed Publications}
\begin{enumerate}
  \pub{Jia Sheng Zhou; Mohd Jaki Bin Mamat; Quan Wen}{2025}{Analyzing the traditional rural cultural heritage of Lingnan region in China as the foundation for protection and development}{International Journal of Conservation Science}{10.36868/ijcs.2025.01.06}{WoS ESCI Q1; co-author.}
  \pub{Quan Wen; Mazran Ismail; Muhammad Hafeez Abdul Nasir}{2024}{Evaluating the residential environment of traditional settlements in northern Jiangxi, China: A multi-dimensional framework}{Journal of Construction in Developing Countries}{10.21315/jcdc.2024.29.s1.11}{WoS ESCI Q4; first author.}
  \pub{Zhang Jing; Hasnanywati Hassan; Zalena Abdul Aziz; Khoo Terh Jing; Wen Quan}{2025}{Exploring the factors for the use of building information modelling (BIM) in historic buildings}{Journal of Construction in Developing Countries}{10.21315/jcdc.2025.30.s2.14}{WoS ESCI; co-author.}
  \pub{Quan Wen; Mazran Ismail; Muhammad Hafeez Abdul Nasir; Jing Zhao; Zhang Jing}{2025}{Assessment and prediction model for community resilience in China based on GANs and random forest regression}{Planning Malaysia}{10.21837/pm.v23i38.1823}{Scopus-indexed; first author.}
  \pub{Quan Wen; Mazran Ismail; Muhammad Hafeez Abdul Nasir; Jing Zhao; Zhang Jing}{2025}{Developing a comprehensive framework for assessing community resilience to public health emergencies in Henan Province: Insights from SciBERT analysis}{Planning Malaysia}{10.21837/pm.v23i38.1822}{Scopus-indexed; first author.}
\end{enumerate}

\section*{Research Projects}
\begin{itemize}
  \item \textbf{Comprehensive Aging-in-Place Renovation Framework for Legacy Residential Communities in Henan Province} (2022). Henan Provincial Social Science Research Fund, Project No. SKL-2022-608. Role: Principal Investigator.
  \item \textbf{Innovative Universal Design Strategies for New Community Development in Henan Province} (2018). Henan Provincial Social Science Research Fund, Project No. SKL-2018-490. First Prize. Role: Principal Investigator.
  \item \textbf{Interior Environment Design Strategies for Hospital Outpatient Departments in the Era of Telemedicine} (2020). Henan Provincial Social Science Research Fund, Project No. SKL-2020-1566. Role: Principal Investigator.
  \item \textbf{Integrated Healthcare-Housing Model for Community-Based Senior Living Facilities in Zhengzhou} (2018). Zhengzhou Federation of Social Sciences Research Project, Project No. JX20190146. Role: Principal Investigator.
  \item \textbf{Adaptive Renovation Strategies for Existing Long-term Care Facilities under COVID-19 Safety Protocols} (2021). Henan Provincial Department of Education Research Fund, Project No. 2021-ZDJH-433. Role: Principal Investigator.
\end{itemize}

\section*{Teaching and Professional Experience}
\begin{itemize}
  \item Lecturer, Henan Industry and Trade Vocational College, China, 2021--2023.
  \item Part-time Lecturer, Henan University of Animal Husbandry and Economy, China, 2019--2021.
  \item Lecturer, Easton International Art College, Zhengzhou University of Light Industry, China, 2017--2021.
  \item Research Associate, Green Building Research Institute, Henan Academy of Building Research, China, 2017.
\end{itemize}

\section*{Professional Service and Certifications}
\begin{itemize}
  \item Journal Reviewer, \textit{SAGE Open}.
  \item Academic Committee Member and Senior Project Review Expert, Zhongren Art Museum.
  \item Higher Education Teaching Qualification Certificate.
  \item Lecturer Certificate, Zhengzhou University of Light Industry.
  \item National Computer Proficiency Examination Certification.
  \item Network Engineering Professional Certification.
\end{itemize}

\section*{Languages}
\begin{itemize}
  \item Chinese Mandarin: Native.
  \item English: Advanced academic and professional proficiency; TOEFL 86/120.
  \item Malay: Conversational proficiency.
\end{itemize}

\section*{Referee}
\textbf{Dr. Mazran bin Ismail}, Senior Lecturer, Architecture, School of Housing, Building and Planning, Universiti Sains Malaysia. Email: \href{mailto:mazran@usm.my}{mazran@usm.my}. Relationship: PhD supervisor.

\end{document}
